% ============================================================
% RESENHA CRÍTICA
% Texto resenhado: FRAGA, Paulo Denisar. Lugar das Humanidades
% na ideia de Universidade crítica. In: FERRAREZI JR., Celso (Org.).
% A identidade docente no Ensino Superior e a universidade brasileira.
% São Paulo; Alfenas: Scortecci; Unifal-MG, 2011. p. 177-195.
%
% Disciplina: Filosofia e Metodologia da Ciência (UNIFAL-MG)
% ============================================================

\documentclass[12pt]{article}

% ---------- Pacotes de codificação e idioma ----------
\usepackage[T1]{fontenc}
\usepackage[utf8]{inputenc}
\usepackage[brazilian]{babel}

% ---------- Margens (3 cm superior/inferior, 2,5 cm laterais) ----------
\usepackage[top=3cm,left=3cm,bottom=2cm,right=2cm]{geometry}

% ---------- Recuo de parágrafo e espaçamento ----------
\usepackage{indentfirst}      % indenta também o primeiro parágrafo de cada seção
\usepackage{setspace}         % controle de espaçamento entre linhas
\setlength{\parindent}{1.25cm}
\setlength{\parskip}{0pt}

% ---------- Pacotes adicionais solicitados ----------
\usepackage{hyperref}
\usepackage{graphicx}
\usepackage{amsmath}
\usepackage{amssymb}

% Links discretos (sem caixas coloridas), conforme sobriedade acadêmica
\hypersetup{
    colorlinks=true,
    linkcolor=black,
    urlcolor=black,
    citecolor=black,
    pdftitle={Resenha - Lugar das Humanidades na ideia de Universidade crítica},
    pdfauthor={Jeann}
}

\begin{document}

% ============================================================
% CABEÇALHO DE IDENTIFICAÇÃO (espaçamento simples, sem recuo)
% ============================================================
\begin{singlespace}
\begin{center}
UNIVERSIDADE FEDERAL DE ALFENAS -- UNIFAL-MG\\
Curso de Ciência da Computação\\
Disciplina: Filosofia e Metodologia da Ciência\\
Docente: José Francisco Lopes Xarão\\
Discente: Jeann Victor Batista\\[6mm]
{\bfseries\Large RESENHA}
\end{center}

\vspace{4mm}
\noindent
FRAGA, Paulo Denisar. Lugar das Humanidades na ideia de Universidade crítica.
\textbf{In}: FERRAREZI JR., Celso (Org.). \textbf{A identidade docente no
Ensino Superior e a universidade brasileira}: uma contribuição da
Universidade Federal de Alfenas ao debate nacional. São Paulo; Alfenas:
Scortecci; Unifal-MG, 2011. p. 177-195.
\end{singlespace}

\vspace{6mm}

% ============================================================
% CORPO DO TEXTO (espaçamento 1,5, conforme solicitado)
% ============================================================
\onehalfspacing

% ---------------------------------------------------------
\section*{Apresentação da obra e do autor}

O texto \emph{Lugar das Humanidades na ideia de Universidade crítica} é de
autoria de Paulo Denisar Vasconcelos Fraga, professor do Instituto de
Ciências Humanas e Letras (ICHL) da Universidade Federal de Alfenas
(UNIFAL-MG), instituição da qual foi também diretor. O autor é licenciado em
Filosofia pela Universidade Federal de Santa Maria e mestre e doutor em
Filosofia pela Universidade Estadual de Campinas (Unicamp), atuando
principalmente nas áreas de Teoria Crítica da Sociedade, Idealismo Alemão e
Filosofia Política, com pesquisas voltadas a autores como Hegel, Marx,
Adorno e Benjamin. É digno de nota que Fraga tenha realizado estágio de
pesquisa na Humboldt-Universität zu Berlin, instituição que herdou o nome
do próprio idealizador do modelo de universidade que serve de referência
central ao capítulo aqui resenhado.

O texto corresponde originalmente a um painel apresentado na mesa-redonda
``A responsabilidade da universidade na formação do sujeito crítico'',
durante o VII Seminário sobre Leitura e Produção no Ensino Superior,
realizado na própria UNIFAL-MG em setembro de 2010. Posteriormente, o
painel foi transformado em capítulo do livro organizado por Celso Ferrarezi
Jr., publicado em 2011 como parte de uma coletânea voltada a discutir a
identidade docente no ensino superior brasileiro a partir da experiência da
UNIFAL-MG.

% ---------------------------------------------------------
\section*{Tema, problema e hipótese}

O tema central do texto é o lugar das Humanidades na formação do sujeito
crítico no interior da Universidade. O problema que orienta toda a
exposição pode ser formulado nos seguintes termos: a responsabilidade de
formar sujeitos críticos depende apenas do aperfeiçoamento de métodos e
didáticas de ensino, ou depende de algo mais estrutural, ligado à própria
concepção de Universidade adotada? Fraga rejeita explicitamente a primeira
alternativa e apresenta sua hipótese norteadora logo na abertura do texto:
a de que ``a efetivação de tal responsabilidade depende primeiramente da
própria concepção de Universidade que tivermos'' (FRAGA, 2011, p. 178), e,
mais especificamente, do papel que nela cumprem as Humanidades -- hipótese
que o autor desenvolve retomando a perspectiva humboldtiana de Universidade
moderna.

% ---------------------------------------------------------
\section*{Síntese do conteúdo}

O texto está organizado em quatro partes. Na primeira, ``Hipótese
norteadora e modo de exposição'', o autor situa o problema acima e anuncia
que irá desenvolvê-lo de modo não sistemático, recorrendo a imagens e
passagens de diferentes pensadores -- método que ele próprio associa ao
gesto benjaminiano de escovar a história a contrapelo (FRAGA, 2011, p.
177-179).

Na segunda parte, ``Passagens sobre Ciência e reflexão humanística'', o
autor reúne cinco exemplos para sustentar a necessidade de mediação entre
as Ciências Empíricas e as Humanidades. Recorre a Franklin Leopoldo e Silva
para discutir por que as Humanidades costumam ser vistas como ``arcaicas''
frente a uma Ciência tida como ``moderna''; retoma a metáfora bachelardiana
da casa, do porão e do sótão -- relida por Gérard Fourez a partir da
distinção entre ``código restrito'' e ``código elaborado'' -- para
contrastar o ``como'' técnico-científico com o ``porquê'' reflexivo próprio
das Humanidades; recupera a disputa entre Demócrito e Epicuro sobre a
natureza do átomo, retomada na tese de doutorado do jovem Marx, para
mostrar como a abertura ao acaso (o \emph{clinamen}) já apontava, no plano
físico, para a exigência de liberdade e de subjetividade crítica; recorre à
biografia de Einstein -- sua correspondência com Freud e seu posterior
arrependimento por ter incentivado o Projeto Manhattan -- para ilustrar os
limites do controle que o próprio cientista tem sobre o uso social de sua
descoberta; e, por fim, aponta o surgimento da Bioética como exemplo de uma
ciência natural que precisou ser cotejada pela reflexão ética sobre seus
limites e finalidades (FRAGA, 2011, p. 179-187).

Na terceira parte, ``A relação entre Ciências e Humanidades no conceito
histórico de Universidade moderna'', o autor recupera o modelo de Wilhelm
von Humboldt, idealizador da Universidade de Berlim, destacando dois
elementos centrais: a unidade indissociável entre pesquisa e ensino
(\emph{Forschung}), que supera a relação tradicional de mera transmissão
entre mestre e discípulo; e a articulação entre ensino, Ciência e
Filosofia, sendo esta última responsável por assegurar a unidade do saber
fragmentado. Fraga cita o próprio Memorando de Humboldt, para quem ``a
Ciência há de ser considerada como algo ainda não de todo encontrado, e que
nunca pode sê-lo, devendo ser buscada ininterruptamente como tal''
(HUMBOLDT, 2008, p. 183 \emph{apud} FRAGA, 2011, p. 189), e recupera ainda
a leitura de Habermas segundo a qual o projeto berlinense se fundava na
unidade entre investigação e ensino, entre ciência e cultura geral e entre
ciência e esclarecimento crítico (FRAGA, 2011, p. 188-191).

Na quarta e última parte, ``Das contradições na cultura e da exigência do
sujeito crítico'', o autor reconhece que as próprias Humanidades não estão
isentas de contradições, encontrando-se também atravessadas pela corrente
positivista que buscam criticar. O texto se encerra parafraseando Sérgio
Paulo Rouanet para afirmar que o ``não lugar'' da Filosofia, da História e
da Literatura na Universidade corresponde, respectivamente, à perda de um
pensamento questionador e relacionante dos saberes, de uma consciência
histórica capaz de perceber o presente como transformável e de um
imaginário capaz de conceber futuros distintos do presente (FRAGA, 2011, p.
192-194).

% ---------------------------------------------------------
\section*{Apreciação crítica}

Do ponto de vista interpretativo, o aspecto mais instigante da obra é a apropriação da metáfora do ``apartamento'', do ``porão'' e do ``sótão'', apresentada por Fraga a partir de Bachelard e Fourez. A ideia de viver em um único plano, aceitando apenas uma perspectiva sobre a realidade, remete diretamente ao romance \emph{1984}, de George Orwell, no qual o controle sobre os indivíduos passa também pela limitação da capacidade de questionar aquilo que lhes é apresentado como verdade. Nesse sentido, a metáfora reforça uma das ideias mais importantes do texto: a Universidade não deve apenas transmitir conhecimentos, mas proporcionar ao indivíduo diferentes perspectivas para a construção da autonomia crítica. O ``sótão'' e o ``porão'', portanto, representam justamente a possibilidade de sair do ``apartamento'' e observar a realidade por outros ângulos, questionando o que, à primeira vista, parece incontestável.

Nota-se, contudo, uma lacuna perceptível: embora o texto critique a supremacia da racionalidade técnico-científica, dialoga-se muito pouco com as próprias Ciências Exatas e Tecnológicas contemporâneas para além de exemplos históricos, como a física de Einstein e o atomismo grego. Escrito em 2010, o trabalho antecede a centralidade que temas como grandes bases de dados, algoritmos de recomendação e inteligência artificial adquiririam na década seguinte -- justamente domínios em que a tensão entre o ``código restrito'' (o desempenho técnico) e o ``código elaborado'' (o sentido e as implicações sociais desse desempenho), descrita a partir de Fourez, tornou-se ainda mais urgente e concreta.

Quanto ao público-alvo, a leitura é especialmente recomendada a estudantes e pesquisadores de licenciaturas, de Pedagogia e de Ciências Humanas em geral, interessados na discussão sobre a função social da Universidade. Ao mesmo tempo, o texto revela-se profundamente relevante para discentes de áreas técnicas e exatas -- como as Engenharias e a Ciência da Computação -- que buscam refletir sobre os limites de uma formação estritamente instrumental e sobre o papel da criticidade em cursos voltados, à primeira vista, apenas ao domínio de ferramentas.

A referida reflexão ressoa de modo bastante direto com a formação acadêmica em Ciência da Computação. Grande parte da matriz curricular da área é organizada em torno do que Fraga, a partir de Fourez, chama de ``código restrito'': o estudante aprende a resolver problemas, otimizar algoritmos e implementar soluções, mas raramente é instigado, no escopo das disciplinas técnicas, a perguntar pelo sentido social dos artefatos que ajudará a produzir. A crítica do autor a uma instituição reduzida à prestação de serviços -- mencionada a partir de Francisco de Oliveira como ``universidade de resultados'' -- não é um problema circunscrito às Humanidades. Tal lógica define, silenciosamente, boa parte da estruturação dos cursos de tecnologia na contemporaneidade, sobretudo diante da pressão do mercado por egressos imediatamente empregáveis. 

Se, como afirma Fraga, o não lugar da Filosofia é o não lugar de um pensamento questionador, conclui-se que, nas graduações de perfil puramente técnico, essa ausência tende a ser a regra, e não a exceção. É justamente por isso que exercícios reflexivos externos ao domínio estritamente técnico possuem imensurável valor: eles obrigam o futuro desenvolvedor a pensar não apenas no ``como'' de um sistema tecnológico, mas, impreterivelmente, no ``porquê'' e ``para quem'' ele é construído.

\newpage % <--- ESTE COMANDO QUEBRA PARA UMA NOVA PÁGINA

% ============================================================
% REFERÊNCIAS
% ============================================================
\section*{Referências}

\begin{singlespace}
\begin{flushleft} % ABNT exige alinhamento exclusivamente à esquerda nas referências

FRAGA, Paulo Denisar. Lugar das Humanidades na ideia de Universidade
crítica. \textbf{In}: FERRAREZI JR., Celso (Org.). \textbf{A identidade
docente no Ensino Superior e a universidade brasileira}: uma contribuição
da Universidade Federal de Alfenas ao debate nacional. São Paulo; Alfenas:
Scortecci; Unifal-MG, 2011. p. 177-195.

\vspace{12pt} % Espaçamento exigido pela ABNT entre uma referência e outra

ORWELL, George. \textbf{1984}. São Paulo: Companhia das Letras, 2009.

\end{flushleft}
\end{singlespace}

\vspace{4mm}

\begin{singlespace}
\section*{Declaração de uso de Inteligência Artificial}
Para a elaboração desta resenha, foram utilizadas as ferramentas
ChatGPT e Claude como instrumentos de apoio à revisão e organização
textual. As ferramentas foram utilizadas para auxiliar na reformulação
de trechos e na organização das ideias, permanecendo sob responsabilidade
do autor a seleção das ideias, a análise crítica, a redação final e a
verificação das informações apresentadas.
\end{singlespace}

\end{document}
